% !TEX root = ../../../main.tex
% 来源：中学数学实验教材·第一册（代数）
% 章节：因式分解与余式定理 / 因式分解
% 原始文件：5.tex，行 502-533
% 类型：tikz
% ---
    \begin{tikzpicture}
 \begin{scope}
\node (A) at (-.6,1.2) {$1$};
\node (B) at (.6,1.2) {$1$};
\node (C) at (-.6,0) {$1$};
\node (D) at (.6,0) {$6$};
\draw[very thick] (A)--(D);
\draw[very thick] (B)--(C);
\node at (0,-.5){$1+6=7$};
\node at (0,-1){（不合适）};
 \end{scope}   
 \begin{scope}[xshift=3cm]
    \node (A) at (-.6,1.2) {$1$};
    \node (B) at (.6,1.2) {$2$};
    \node (C) at (-.6,0) {$1$};
    \node (D) at (.6,0) {$3$};
    \draw[very thick] (A)--(D);
    \draw[very thick] (B)--(C);
    \node at (0,-.5){$2+3=5$};
    \node at (0,-1){（不合适）};
     \end{scope} 
     \begin{scope}[xshift=6cm]
        \node (A) at (-.6,1.2) {$1$};
        \node (B) at (.6,1.2) {$-2$};
        \node (C) at (-.6,0) {$1$};
        \node (D) at (.6,0) {$-3$};
        \draw[very thick] (A)--(D);
        \draw[very thick] (B)--(C);
        \node at (0,-.5){$-2+(-3)=-5$};
        \node at (0,-1){（合适）};
         \end{scope} 
\end{tikzpicture}
